\begin{abstract}
Machine learning researchers often discover computational feasibility constraints only after full implementation, wasting hours to days on non-viable hypotheses. Existing approaches fail to provide early-stop protocols: ablation studies test hypothesis \textit{variations} assuming viability rather than assessing viability itself, Big-O analysis misses constant factors and hardware specifics, and expert intuition remains informal with untested accuracy. We introduce Pilot-Driven Viability Gates, a framework that treats viability assessment as incremental empirical validation across sample scales (10 $\rightarrow$ 100 $\rightarrow$ full samples) with Bayesian posterior refinement. Our key insight: computational overhead scales predictably from micro-pilot to full dataset, enabling accurate early prediction. In validation on a corpus of 32 synthetic ML hypotheses under a 10\% overhead threshold, Gate 1 (10-sample micro-pilot) achieved 93.3\% accuracy (28 of 30 correct) for identifying non-viable hypotheses, exceeding our 80\% target and significantly outperforming random baseline (binomial p=4.34e-07). The framework exhibited perfect linear correlation (r=1.000, p<0.0001) between micro-pilot and full-scale overhead, with Bayesian updates reducing prediction error by 40.91\% (p=0.0003). While validated on synthetic data---establishing proof-of-concept for framework mechanics---external validity requires real-world corpus validation. Our work provides the first formalized feasibility-first methodology for early identification of computationally infeasible hypotheses, addressing a gap in ML research workflows where systematic viability gates are absent.
\end{abstract}
